\documentclass[10pt,a4paper]{article}
\usepackage{fontspec,amsmath,amssymb,graphicx}
\usepackage[margin=18mm]{geometry}
\setmainfont{Times New Roman}
\setlength{\parindent}{0pt}
\setlength{\parskip}{5pt}
\allowdisplaybreaks
\sloppy
\emergencystretch=2em
\begin{document}
\section*{BT\_2\_3\_12 — Lời giải bổ sung}
\subsection*{Kiến thức và trình tự giải}
\textbf{Mục tiêu.} Giải phương trình hyperbolic sau khi đã tìm hai tọa độ đặc trưng độc lập.

\textbf{Bước 1 — Phân loại, tìm đặc trưng và thay đạo hàm.} Làm đủ các bước của phép đưa về chính tắc. Viết cả phép đổi ngược để biểu diễn vế phải $\widetilde f(\xi,\eta)=f(x(\xi,\eta),y(\xi,\eta))$.

\textbf{Bước 2 — Tích phân phương trình chính tắc.} Với $v_{\xi\eta}=g$, giữ biến kia cố định khi tích phân:


\[
v=F(\xi)+G(\eta)+\int_{\xi_0}^{\xi}\int_{\eta_0}^{\eta}g(a,b)\,db\,da.
\]


Nếu tìm được một nghiệm riêng $P(x,y)$ dễ hơn, viết $u=P+F(\xi)+G(\eta)$. Với $v_{\xi\eta}+\alpha v_\eta=0$, đặt $w=v_\eta$, giải $w_\xi+\alpha w=0$ rồi tích phân theo $\eta$; thu được $v=F(\xi)+e^{-\alpha\xi}G(\eta)$.

\textbf{Bước 3 — Đổi cả dữ liệu.} Thế tọa độ đường dữ liệu vào $\xi,\eta$, dùng quy tắc dây chuyền cho đạo hàm được cho. Trên đường không đặc trưng, hai dữ liệu thường xác định hai hàm $F,G$. Trên đường đặc trưng, có biến không thay đổi, nên dữ liệu chỉ xác định giá trị hoặc đạo hàm của một hàm tại một điểm; phải tìm điều kiện tương thích và nêu phần tùy ý còn lại.

\textbf{Bước 4 — Lập đẳng thức hình bình hành.} Với $v=F(\xi)+G(\eta)$, bốn đỉnh theo hai họ đặc trưng thỏa $v_{11}+v_{22}=v_{12}+v_{21}$. Nếu có số hạng mũ, trước hết nhân trọng số để được tổng hai hàm một biến. Với vế phải $g$, hiệu bốn đỉnh bằng tích phân kép của $g$ trên hình chữ nhật trong tọa độ mới.

\textbf{Bước 5 — Xác định miền được dữ liệu quyết định.} Hai đặc trưng đi qua điểm đang xét phải gặp đoạn dữ liệu phù hợp. Với phương trình sóng $v_{tt}=v_{zz}$ trên $z>0$, dữ liệu lúc $t=0$ xác định duy nhất miền $z>t$; ở $0<z<t$, cần thêm dữ liệu biên $z=0$. Với đặc trưng cong, dùng miền trong các tọa độ mới rồi đổi lại.

\textbf{Bước 6 — Kiểm tra.} Thế nghiệm vào cả PDE và từng điều kiện. Nếu đề yêu cầu hai nghiệm, chỉ rõ hai lựa chọn khác nhau cho hàm tùy ý và chứng minh phần chênh lệch thỏa phương trình thuần nhất, đồng thời triệt tiêu trên đường dữ liệu.
\subsection*{Lời giải áp dụng cho bài}
\textbf{(a) Phân loại và đổi biến.} Trong các biến $(x,t)$, $A=x$, $B=0$, $C=-x^5$, nên $\delta=x^6>0$ ở $x>0$. Chọn $z=x^3/3$, $u(x,t)=v(z,t)$, $z_x=x^2>0$. Tính
\[u_x=x^2v_z,\qquad u_{xx}=2xv_z+x^4v_{zz}.\]
Hai hạng cấp một triệt tiêu: $xu_{xx}-2u_x=x^5v_{zz}$. PDE trở thành $v_{tt}=v_{zz}$; dạng đặc trưng $(\xi,\eta)=(z+t,z-t)$ cho $v_{\xi\eta}=0$.

\textbf{(b) Tổng quát.} Tích phân đạo hàm hỗn hợp bằng $0$ cho $u=F(x^3/3+t)+G(x^3/3-t)$. Hai họ đặc trưng là $x^3/3\pm t=\mathrm{const}$.

\textbf{(c) Dữ liệu ban đầu.} Đặt $\Phi(z)=\varphi((3z)^{1/3})$, $\Psi(z)=\psi((3z)^{1/3})$, $z>0$. Tại $t=0$, $F+G=\Phi$, $F'-G'=\Psi$. Lấy đạo hàm phương trình đầu rồi giải hệ: $F'=(\Phi'+\Psi)/2$, $G'=(\Phi'-\Psi)/2$. Do đó công thức d'Alembert là
\[\boxed{u=\frac{\Phi(z+t)+\Phi(z-t)}2+\frac12\int_{z-t}^{z+t}\Psi(s)ds,\quad z=x^3/3>t.}\]
Miền duy nhất do dữ liệu ban đầu là $x>0$, $0<t<x^3/3$. Trong $0<z<t$, một đặc trưng gặp biên $z=0$ trước khi gặp trục $t=0$, nên cần thêm dữ liệu biên. Công thức kiểm tra $u(x,0)=\varphi(x)$, $u_t(x,0)=\psi(x)$ và $v_{tt}=v_{zz}$.

\textbf{(iv) Chọn điều khiển để $u(x,9)=0$ với $0\le x\le3$.} Đoạn đích tương ứng $0\le z\le9$. Giả sử $\Phi,\Psi$ có các giới hạn và độ trơn cần thiết tại $z=0$. Với $0\le z\le t$, dữ liệu biên $v(0,\tau)=h(\tau)$ cho
\[v(z,t)=h(t-z)+\frac{\Phi(t+z)-\Phi(t-z)}2+\frac12\int_{t-z}^{t+z}\Psi(s)ds.\]
Cách suy ra: $F(t+z)+G(z-t)$ và $h(t-z)=F(t-z)+G(z-t)$, nên hiệu bằng $F(t+z)-F(t-z)$; dùng $F'=(\Phi'+\Psi)/2$. Tại $t=9$, đặt $\tau=9-z\in[0,9]$, yêu cầu bằng $0$ cho
\[\boxed{h(\tau)=-\frac{\Phi(18-\tau)-\Phi(\tau)}2-\frac12\int_\tau^{18-\tau}\Psi(s)ds,\quad0\le\tau\le9.}\]
Đây là điều khiển duy nhất trên khoảng thời gian cần dùng; sau $9$ có thể kéo dài đủ trơn. Tại góc $(z,t)=(0,0)$, điều kiện giá trị $h(0)=\Phi(0)$ tương đương
\[\Phi(18)+\Phi(0)+\int_0^{18}\Psi(s)ds=0.\]
Muốn nghiệm $C^2$ qua đặc trưng $z=t$, còn cần $h'(0)=\Psi(0)$ và $h''(0)=\Phi''(0)$. Thay công thức $h$ cho hai điều kiện cụ thể:
\[\Phi'(18)+\Phi'(0)+\Psi(18)-\Psi(0)=0,\qquad
\Phi''(18)+\Phi''(0)+\Psi'(18)-\Psi'(0)=0.\]
Nếu dữ liệu không thỏa điều kiện góc tương ứng, không có nghiệm cổ điển với điều khiển này và đích đã cho trong lớp trơn ấy; không bỏ qua điểm $z=t$ khi kết luận. Với các điều kiện đúng, thế $\tau=9-z$ vào công thức vùng biên cho $v(z,9)=0$ trên toàn đoạn đích.

\textbf{Chi tiết tích phân hệ dữ liệu đầu.} Chọn mốc $z_0>0$ và $F(z_0)=c$. Từ $F'=(\Phi'+\Psi)/2$,
\[F(s)=c+\frac{\Phi(s)-\Phi(z_0)}2+\frac12\int_{z_0}^{s}\Psi(r)dr,\quad G(s)=\Phi(s)-F(s).\]
Thế $F(z+t)+G(z-t)=F(z+t)-F(z-t)+\Phi(z-t)$; hai hằng số mất, các tích phân gộp thành $\int_{z-t}^{z+t}\Psi$, cho đúng d'Alembert với $z-t>0$.\end{document}